% Matching
\question\MatchQuestion{Diphycercal}{Type of tail found in coelacanths.}
\vspace{0.5\baselineskip}
%\vspace{0.4\baselineskip}


\question[15]
Explain the difference between shoaling and schooling.  Explain two possible adaptive reasons for schooling behavior.

	\begin{solution}
	\end{solution}

\vspace*{\stretch{4}}

\question[11]
Using the provided phylogeny of living fishes to fill in the blanks with the appropriate systematic group and hierarchical level (class, subclass, etc.).  Each lettered blank below corresponds to the lettered branch on the phylogeny. ($1/2$ point per blank). Spelling counts throughout.\vspace{\baselineskip}

\begin{tabular}{@{}lll@{}}
	&	\multicolumn{1}{c}{Group Name}		& 	\multicolumn{1}{c}{Classification Level} \\[1.5em]
A.	& 	\rule{2.75in}{0.4pt}	&	\rule{2.75in}{0.4pt} \\[1.5em]
B.	& 	\rule{2.75in}{0.4pt}	&	\rule{2.75in}{0.4pt} \\[1.5em]
C.	& 	\rule{2.75in}{0.4pt}	&	\rule{2.75in}{0.4pt} \\[1.5em]
D.	& 	\rule{2.75in}{0.4pt}	&	\rule{2.75in}{0.4pt} \\[1.5em]
E.	& 	\rule{2.75in}{0.4pt}	&	\rule{2.75in}{0.4pt} \\[1.5em]
F.	& 	\rule{2.75in}{0.4pt}	&	\rule{2.75in}{0.4pt} \\[1.5em]
G.	& 	\rule{2.75in}{0.4pt}	&	\rule{2.75in}{0.4pt} \\[1.5em]
H.	& 	\rule{2.75in}{0.4pt}	&	\rule{2.75in}{0.4pt} \\[1.5em]
I.	& 	\rule{2.75in}{0.4pt}	&	\rule{2.75in}{0.4pt} \\[1.5em]
J.	& 	\rule{2.75in}{0.4pt}	&	\rule{2.75in}{0.4pt} \\[1.5em]
K.	& 	\rule{2.75in}{0.4pt}	&	\rule{2.75in}{0.4pt} \\[1.5em]
\end{tabular}

%Grad Student Questions
\textbf{(GSH)} Name the type of soft fin ray that is diagnostic of the Coelacanthimorpha. How does this ray differ from that same type of soft fin rays found in the Actinopterygii?

%Extra Credit questions
\question[2]
\textbf{Extra Credit:} You discover a small fossil fish that clearly has jaws, heavy bony plating, a dorsoventrally flattened body, pectoral fins (no pelvic fins), and a heterocercal tail.  This fossil fish likely belongs to what class?

